%==================================================================================================
% LaTeX paper template - use as a starting point for structuring a research paper.
%
% Written by Colin Perkins (https://csperkins.org/)
% 2002-2026
%
% To the extent possible under law, the author(s) have dedicated all copyright and
% related and neighbouring rights to this software to the public domain worldwide.
% This template is distributed without any warranty.
%
% You should have received a copy of the CC0 Public Domain Dedication along with
% this software. If not, see <http://creativecommons.org/publicdomain/zero/1.0/>.
%
%   *************************************************************************
%   *** If you use this example.tex file as the basis for your paper, you ***
%   *** MUST remove this header block and the above public domain grant.  ***
%   *************************************************************************

%==================================================================================================
% General advice on technical writing:
%  - George Gopen and Judith Swan, "The Science of Scientific Writing",
%    American Scientist, Nov/Dec 1990.
%    http://www.americanscientist.org/issues/num2/the-science-of-scientific-writing/1
%  - Stephen Pinker, "The Sense of Style: The Thinking Person's Guide to
%    Writing in the 21st Century", Penguin, Sept 2014. ISBN 0525427929.
%
% The paper writing advice in the comments is derived from talks and articles
% by Simon Peyton-Jones, Jim Kurose, Henning Schulzrinne, and Jim Bednar:
%  - http://research.microsoft.com/~simonpj/papers/giving-a-talk/giving-a-talk.htm
%  - http://research.microsoft.com/~simonpj/papers/giving-a-talk/writing-a-paper-slides.pdf
%  - http://gaia.cs.umass.edu/kurose/talks/top_10_tips_for_writing_a_paper.ppt
%  - http://www-net.cs.umass.edu/kurose/writing/intro-style.html
%  - http://www.cs.columbia.edu/~hgs/etc/writing-style.html
%  - http://homepages.inf.ed.ac.uk/jbednar/writingtips.html
%  - http://www.gabbay.org.uk/blog/paper-writing.html
%  - http://homes.cs.washington.edu/~mernst/advice/write-technical-paper.html
%  - http://dx.doi.org/10.1371/journal.pcbi.1005619
%
% LaTeX usage notes:
%  - http://www.read.seas.harvard.edu/~kohler/latex.html

%==================================================================================================
% Recommended preamble for a paper using the ACM template. This will need
% to be adapted for use with other templates: follow the guidance for the
% venue to which you're submitting.
%
% The \DocumentMetadata line enables the accessibility features in LaTeX
% PDF generation. It requires you have TeX Live 2026 installed and build
% using lualatex rather than pdflatex.
%
% The \usepackage lines install commonly used extensions.These should be
% available in most LaTeX installations. Documentation can be found at
% https://www.ctan.org/. Note that the order in which packages are loaded
% is significant.
%
% The \synctex=1 command tells LaTeX to embed information about the source
% file into the PDF. This gives PDF viewers the ability to tell the text
% editors to go to the matching location in the document when the PDF is
% clicked.

\DocumentMetadata{tagging=on,tagging-setup={math/setup=mathml-SE},pdfstandard=ua-2,lang=en-GB}
\documentclass[10pt,sigconf]{acmart}
\usepackage[l2tabu,orthodox]{nag}
\usepackage[british]{babel}
\usepackage{microtype}
\usepackage{amsmath}
% These are part of acmart so don't need to be included here:
%   \usepackage{newtxtext}
%   \usepackage{newtxmath}
\usepackage{subcaption}     % Incompatible with tagging as on 2026-04-17
\usepackage{booktabs}
\usepackage{upquote}
\usepackage{graphicx}
\usepackage{url}
\usepackage{listings}       % Incompatible with tagging as on 2026-04-17
\usepackage{float}
\usepackage{hyperref}
\usepackage{algorithm}      % Incompatible with tagging as on 2026-04-17
\usepackage{color}

\uchyph=0   % Prevent hyphenation of all upper case words:
\synctex=1  % Enable synctex extension

% Define a simple \todo{...} macro:
\newcommand{\todo}[1]{\textbf{\textcolor{red}{To do: #1}}}

%==================================================================================================
\begin{document}
\title{Title goes here}

\author{Colin Perkins}
\orcid{0000-0002-3404-8964}
\affiliation{
  \institution{University of Glasgow}
  \streetaddress{School of Computing Science}
  \city{Glasgow}
  \postcode{G12 8QQ}
  \country{UK}
}
\email{csp@csperkins.org}

\acmYear{2018}
\copyrightyear{2018}
\setcopyright{acmcopyright}
\acmConference{CoNEXT '18}{December 4--7, 2018}{Heraklion/Crete, Greece}
\acmPrice{15.00}
\acmDOI{10.1145/3284850.3284856}
\acmISBN{978-1-4503-6082-1/18/12}

%==================================================================================================
\begin{abstract}
  % Four sentences:
  %  - State the problem
  %  - Say why it's an interesting problem
  %  - Say what your solution achieves
  %  - Say what follows from your solution

  The abstract goes here.

\end{abstract}
\maketitle
%==================================================================================================
\section{Introduction}

% A good paper introduction is fairly formulaic. If you follow a simple set
% of rules, you can write a very good introduction. The following outline can
% be varied. For example, you can use two paragraphs instead of one, or you
% can place more emphasis on one aspect of the intro than another. But in all
% cases, all of the points below need to be covered in an introduction, and
% in most papers, you don't need to cover anything more in an introduction.
%
% Paragraph 1: Motivation. At a high level, what is the problem area you
% are working in and why is it important? It is important to set the larger
% context here. Why is the problem of interest and importance to the larger
% community?



% Paragraph 2: What is the specific problem considered in this paper? This
% paragraph narrows down the topic area of the paper. In the first
% paragraph you have established general context and importance. Here you
% establish specific context and background.



% Paragraph 3: "In this paper, we show that...". This is the key paragraph
% in the introduction - you summarize, in one paragraph, what are the main
% contributions of your paper, given the context established in paragraphs
% 1 and 2. What's the general approach taken? Why are the specific results
% significant? The story is not what you did, but rather:
%  - what you show, new ideas, new insights
%  - why interesting, important?
% State your contributions: these drive the entire paper.  Contributions
% should be refutable claims, not vague generic statements.

In this paper, we ...

% Paragraph 4: What are the differences between your work, and what others
% have done? Keep this at a high level, as you can refer to future sections
% where specific details and differences will be given, but it is important
% for the reader to know what is new about this work compared to other work
% in the area.



% Paragraph 5: "We structure the remainder of this paper as follows." Give
% the reader a road-map for the rest of the paper. Try to avoid redundant
% phrasing, "In Section 2, In section 3, ..., In Section 4, ... ", etc.

We structure the remainder of this paper as follows.

%==================================================================================================
% Concentrate single-mindedly on a narrative that:
%  - Describes the problem, and why it's interesting
%  - Describes your idea
%  - Defends your idea, showing how it solves the problem, and filling out
%    the details
% On the way, cite relevant work in passing, but defer discussion to the
% end.
%
% Introduce the problem, and your idea, using examples, and only then
% present the general case. Explain the idea as if your were speaking to
% someone using a whiteboard. Conveying the intuition is primary; details
% follow. Write in a top down manner: state broad themes and ideas first,
% then go into details.
%
% The introduction makes claims. The body of the paper provides evidence
% to support each claim. Check each claim in the introduction, identify
% the evidence, and forward-reference it from the claim.
%==================================================================================================
\section{}

% The problem



%==================================================================================================
\section{}

% My idea



%==================================================================================================
\section{}

% The details

% Describe results carefully:
%  - clearly state assumptions
%  - give enough information to allow the reader to recreate the results
%  - ensure results are representative; statistically meaningful, etc.
%  - don't overstate results
%  - equally, don't understate them: consider the broader implications



%==================================================================================================
\section{Related Work}

% This should come near the end, and focussing on discussing how your work
% relates to that of others. Any relevant related work should have been
% cited already, so this is not a list of related work, it's a discussion
% of how that work relates.
%
% Why not put related work after the introduction? 1) because describing
% alternative approaches gets between the reader and your idea; and 2)
% because the reader knows nothing about the problem yet, so your
% (carefully trimmed) description of various technical trade-offs is
% absolutely incomprehensible.
%
% When writing the related work:
%  - Give credit to others where it's due; this doesn't diminish the
%    credit you get from your paper.
%  - Acknowledge weaknesses in your approach.
%  - Ensure related work is accurate and up-to-date



%==================================================================================================
\section{Conclusions}



%==================================================================================================
\section{Acknowledgements}

% Acknowledge funding sources.

%==================================================================================================
% Set the bibliography style. Choose one of the following, depending on the
% document class being used:
%
%   \bibliographystyle{abbrv}                 When using article class
%   \bibliographystyle{IEEEtran}              When using IEEE style
%   \bibliographystyle{ACM-Reference-Format}  When using ACM style

\bibliographystyle{ACM-Reference-Format}

% Load the bibliography file(s) for this paper:
\bibliography{example}

%==================================================================================================
\end{document}
% vim: set ts=2 sw=2 tw=75 et ai:
